\section{How a Deployed System Acquires Its Values}
\label{sec:4.2}
A system deployed today acquired what it will and will not do at several distinct stages, run in a fixed order, by different people, using different instruments. The order is the part that matters most and the part a list of techniques hides. It determines who can change what, at what cost, and with how much warning to anybody else.

The stages run like this. A model is pretrained on a very large corpus and then tuned on demonstrations of the behavior somebody wants. It is tuned again on comparisons between candidate outputs, which is where most of what a user would call its manners and its refusals are set. Some of that feedback is now generated by models rather than people, against a written set of principles. Supervision can be aimed at the steps of a reasoning process rather than at the answer it arrives at. Where the system's competence exceeds any available rater's, the oversight itself becomes a research problem. Interpretability and evaluation have moved from commentary on a finished model to instruments applied during its construction. And at the end, at the moment of deployment, behavior is set again by an instruction prepended to every request and by the tools the system is permitted to call.

Each of those is taken twice below, on the test above: what the stage does, and whether what it does could produce a commitment that survives its own teacher.

The answer comes out the same at every stage, and for the same reason each time. Every stage is a place where some party decides what the system will value, and every one of those decisions is revisable by whoever holds that stage. What varies down the list is how legible the decision is, how expensive it is to reverse, and how many people would have to be involved in reversing it. Those differences are real, and they are differences of degree in a property section~\ref{sec:3.1} shows is not enough at any degree.

That is a location rather than a dismissal. In section~\ref{sec:2.3}'s vocabulary, this pipeline is machinery for producing moral performance; its better stages reach moral competence, which is a genuine achievement and the right target for most of what gets built. What none of them reaches is the third thing, and the reason is the same every time — what a system learns from a party is what that party can make it unlearn.
